\begin{table}[htbp]
\centering
\footnotesize
\caption{Análise de Robustez e Sensibilidade da PIP — Desfecho: Letalidade}
\label{tab:sens_letalidade}
\begin{tabular}{lccccr}
\toprule
\textbf{Covariável} & \textbf{Principal} & \textbf{Ampliado} & \textbf{Em Nível} & \textbf{Sem Fortaleza} & \textbf{$N \ge 0{,}50$} \\
\midrule
Município integrante de região metropolitana & 0,781 & 0,732 & 0,706 & 0,649 & 4 \\
Proporção de mulheres eleitas & 0,744 & 0,697 & 0,686 & 0,626 & 4 \\
Cobertura média de imunização em menores de um ano & 0,635 & 0,591 & 0,706 & 0,534 & 4 \\
Alinhamento partidário do prefeito com o governo federal & 0,616 & 0,559 & 0,580 & 0,528 & 4 \\
IDEB dos anos finais do ensino fundamental & 0,566 & 0,515 & 0,543 & 0,475 & 3 \\
IVS Capital Humano & 0,564 & 0,490 & 0,578 & 0,468 & 2 \\
Leitos vinculados ao SUS & 0,564 & 0,500 & 0,374 & 0,456 & 2 \\
Proporção da população em extrema pobreza & 0,547 & 0,480 & 0,468 & 0,451 & 1 \\
Internações por doenças do aparelho respiratório & 0,536 & 0,479 & 0,568 & 0,437 & 2 \\
Unidades de saúde vinculadas ao SUS & 0,534 & 0,487 & 0,446 & 0,435 & 1 \\
População beneficiária do Programa Bolsa Família & 0,528 & 0,476 & 0,443 & 0,440 & 1 \\
Área territorial & 0,482 & 0,418 & 0,437 & 0,392 & 0 \\
IDEB do 3º ano do ensino médio & 0,481 & 0,431 & 0,432 & 0,406 & 0 \\
Internações por doenças do aparelho circulatório & 0,475 & 0,422 & 0,479 & 0,398 & 0 \\
Proporção da população em empregos formais & 0,468 & 0,414 & 0,491 & 0,386 & 0 \\
Taxa de homicídios & 0,456 & 0,403 & 0,401 & 0,371 & 0 \\
Índice de Desenvolvimento Municipal & 0,456 & 0,417 & 0,423 & 0,380 & 0 \\
População estimada & 0,454 & 0,404 & 0,517 & 0,341 & 1 \\
Índice de Desenvolvimento Humano Municipal & 0,453 & 0,402 & 0,467 & 0,368 & 0 \\
IVS Infraestrutura Urbana & 0,450 & 0,400 & 0,375 & 0,375 & 0 \\
Densidade demográfica & 0,446 & 0,399 & 0,485 & 0,328 & 0 \\
IDEB dos anos iniciais do ensino fundamental & 0,431 & 0,385 & 0,405 & 0,360 & 0 \\
Consumo de energia elétrica por 100 mil habitantes & 0,431 & 0,386 & 0,401 & 0,357 & 0 \\
Profissionais de saúde vinculados ao SUS & 0,426 & 0,378 & 0,464 & 0,351 & 0 \\
Produto Interno Bruto per capita & 0,423 & 0,374 & 0,386 & 0,338 & 0 \\
Proporção da população rural & 0,417 & 0,376 & 0,405 & 0,338 & 0 \\
Proporção da população com 60 anos ou mais & 0,416 & 0,371 & 0,382 & 0,332 & 0 \\
IVS Renda e Trabalho & 0,413 & 0,368 & 0,394 & 0,341 & 0 \\
Proporção de votos válidos do prefeito eleito & 0,412 & 0,366 & 0,388 & 0,329 & 0 \\
Despesa municipal com saúde e saneamento & 0,410 & 0,368 & 0,415 & 0,340 & 0 \\
Município integrante do semiárido & 0,408 & 0,360 & 0,369 & 0,334 & 0 \\
Despesa municipal com educação e cultura & 0,406 & 0,361 & 0,388 & 0,328 & 0 \\
Internações por diabetes mellitus & 0,405 & 0,359 & 0,449 & 0,328 & 0 \\
Internações por asma & 0,405 & 0,361 & 0,379 & 0,329 & 0 \\
Cobertura urbana de abastecimento de água & 0,401 & 0,351 & 0,370 & 0,326 & 0 \\
Auxílio emergencial & - & 0,390 & - & - & 0 \\
Transferências federais para combate à COVID-19 & - & 0,379 & - & - & 0 \\
\bottomrule
\end{tabular}
\begin{tablenotes}\footnotesize
\item \textit{Nota:} Colunas representam a Probabilidade de Inclusão Posterior (PIP) em cada especificação metodológica.
\end{tablenotes}
\end{table}